\begin{ex}
Trong hộp chỉ có 5 viên bi màu đỏ. Nếu nhắm mắt lấy ra 1 viên bi, sự kiện 'Lấy được viên bi màu đỏ' là chắc chắn, có thể hay không thể? Quy ước: Chắc chắn = 1, Có thể = 2, Không thể = 3.
\shortans{1}
\loigiai{
Vì toàn bi đỏ nên chắc chắn lấy được bi đỏ. Đáp án là 1.
}
\end{ex}

\begin{ex}
Trong hộp đó chỉ có bi đỏ, sự kiện 'Lấy được viên bi màu xanh' xảy ra như thế nào? (Quy ước: Chắc chắn = 1, Có thể = 2, Không thể = 3).
\shortans{3}
\loigiai{
Không có bi xanh nên không thể xảy ra. Đáp án là 3.
}
\end{ex}

\begin{ex}
Trong túi có 3 viên bi đỏ và 3 viên bi xanh. Nhắm mắt lấy 1 viên bi, sự kiện 'Lấy được viên bi màu đỏ' là: Chắc chắn = 1, Có thể = 2, Không thể = 3.
\shortans{2}
\loigiai{
Có thể lấy được bi đỏ hoặc bi xanh. Đáp án là 2.
}
\end{ex}

\begin{ex}
Gieo một con xúc xắc 6 mặt (từ 1 chấm đến 6 chấm). Sự kiện 'Xuất hiện mặt có 7 chấm' xảy ra như thế nào? (Chắc chắn = 1, Có thể = 2, Không thể = 3).
\shortans{3}
\loigiai{
Xúc xắc chỉ có từ 1 đến 6 chấm nên không thể có mặt 7 chấm. Đáp án là 3.
}
\end{ex}

\begin{ex}
Gieo một con xúc xắc 6 mặt. Sự kiện 'Xuất hiện mặt có số chấm nhỏ hơn 7' là: (Chắc chắn = 1, Có thể = 2, Không thể = 3).
\shortans{1}
\loigiai{
Mọi mặt đều từ 1 đến 6 (nhỏ hơn 7) nên chắc chắn xảy ra. Đáp án là 1.
}
\end{ex}

\begin{ex}
Gieo một con xúc xắc. Sự kiện 'Xuất hiện mặt 6 chấm' là: (Chắc chắn = 1, Có thể = 2, Không thể = 3).
\shortans{2}
\loigiai{
Có thể xuất hiện mặt 6 chấm. Đáp án là 2.
}
\end{ex}

\begin{ex}
Ngày mai Mặt Trời mọc ở hướng Đông. Sự kiện này là: (Chắc chắn = 1, Có thể = 2, Không thể = 3).
\shortans{1}
\loigiai{
Mặt Trời luôn mọc ở hướng Đông nên chắc chắn xảy ra. Đáp án là 1.
}
\end{ex}

\begin{ex}
Ngày mai Mặt Trời lặn ở hướng Đông. Sự kiện này là: (Chắc chắn = 1, Có thể = 2, Không thể = 3).
\shortans{3}
\loigiai{
Mặt Trời lặn ở hướng Tây nên không thể xảy ra. Đáp án là 3.
}
\end{ex}

\begin{ex}
Một tuần lễ có 7 ngày. Sự kiện này là: (Chắc chắn = 1, Có thể = 2, Không thể = 3).
\shortans{1}
\loigiai{
Đúng quy luật tự nhiên nên chắc chắn xảy ra. Đáp án là 1.
}
\end{ex}

\begin{ex}
Trong đĩa có 5 quả cam và 4 quả táo. Nam lấy ra 6 quả. Sự kiện 'Trong số quả lấy ra có ít nhất 1 quả cam' là: (Chắc chắn = 1, Có thể = 2, Không thể = 3).
\shortans{1}
\loigiai{
Vì chỉ có 4 quả táo nên lấy 6 quả thì chắc chắn phải có quả cam. Đáp án là 1.
}
\end{ex}

\begin{ex}
Trong hộp có 10 chiếc bút chì màu đen. Lấy ra 2 chiếc bút, cả 2 chiếc đều màu đen là: (Chắc chắn = 1, Có thể = 2, Không thể = 3).
\shortans{1}
\loigiai{
Chắc chắn vì toàn bộ bút trong hộp đều màu đen. Đáp án là 1.
}
\end{ex}

\begin{ex}
Bé Nam 8 tuổi cao 3 mét. Sự kiện này là: (Chắc chắn = 1, Có thể = 2, Không thể = 3).
\shortans{3}
\loigiai{
Không thể xảy ra đối với con người. Đáp án là 3.
}
\end{ex}

\begin{ex}
Hôm nay là Thứ Bảy, ngày mai là Chủ Nhật. Sự kiện này là: (Chắc chắn = 1, Có thể = 2, Không thể = 3).
\shortans{1}
\loigiai{
Chắc chắn xảy ra theo thứ tự các ngày trong tuần. Đáp án là 1.
}
\end{ex}

\begin{ex}
Tung một đồng xu có hai mặt Sấp và Ngửa. Sự kiện 'Đồng xu rơi xuống xuất hiện mặt Sấp' là: (Chắc chắn = 1, Có thể = 2, Không thể = 3).
\shortans{2}
\loigiai{
Có thể xuất hiện mặt Sấp hoặc mặt Ngửa. Đáp án là 2.
}
\end{ex}

\begin{ex}
Trong túi có 5 quả bóng xanh và 5 quả bóng vàng. Lấy ra 6 quả bóng, sự kiện 'Có ít nhất 1 quả bóng xanh' là: (Chắc chắn = 1, Có thể = 2, Không thể = 3).
\shortans{1}
\loigiai{
Chắc chắn vì bóng vàng chỉ có 5 quả, lấy 6 quả chắc chắn có bóng xanh. Đáp án là 1.
}
\end{ex}

\begin{ex}
Lấy ra từ túi đó 1 quả bóng màu tím. Sự kiện này là: (Chắc chắn = 1, Có thể = 2, Không thể = 3).
\shortans{3}
\loigiai{
Trong túi không có bóng tím nên không thể. Đáp án là 3.
}
\end{ex}

\begin{ex}
Tháng Hai có 35 ngày. Sự kiện này là: (Chắc chắn = 1, Có thể = 2, Không thể = 3).
\shortans{3}
\loigiai{
Tháng Hai nhiều nhất chỉ có 29 ngày nên không thể có 35 ngày. Đáp án là 3.
}
\end{ex}

\begin{ex}
Một năm có 12 tháng. Sự kiện này là: (Chắc chắn = 1, Có thể = 2, Không thể = 3).
\shortans{1}
\loigiai{
Chắc chắn đúng. Đáp án là 1.
}
\end{ex}

\begin{ex}
Số 24 là số chẵn. Sự kiện này là: (Chắc chắn = 1, Có thể = 2, Không thể = 3).
\shortans{1}
\loigiai{
Chắc chắn đúng. Đáp án là 1.
}
\end{ex}

\begin{ex}
Tổng của hai số lẻ là một số lẻ. Sự kiện này là: (Chắc chắn = 1, Có thể = 2, Không thể = 3).
\shortans{3}
\loigiai{
Tổng hai số lẻ luôn là số chẵn nên không thể là số lẻ. Đáp án là 3.
}
\end{ex}

\begin{ex}
Chọn ngẫu nhiên một bạn học sinh trong lớp 2A toàn nữ. Bạn được chọn là học sinh nữ: (Chắc chắn = 1, Có thể = 2, Không thể = 3).
\shortans{1}
\loigiai{
Toàn nữ nên chắc chắn là học sinh nữ. Đáp án là 1.
}
\end{ex}

\begin{ex}
Bạn được chọn là học sinh nam: (Chắc chắn = 1, Có thể = 2, Không thể = 3).
\shortans{3}
\loigiai{
Không có bạn nam nào nên không thể. Đáp án là 3.
}
\end{ex}

\begin{ex}
Thời tiết ngày mai có mưa: (Chắc chắn = 1, Có thể = 2, Không thể = 3).
\shortans{2}
\loigiai{
Có thể mưa hoặc không mưa. Đáp án là 2.
}
\end{ex}

\begin{ex}
Có bao nhiêu sự kiện 'Chắc chắn' trong 3 sự kiện sau: (A) Cá sống dưới nước; (B) Con người bay bằng cánh tay; (C) 1 ngày có 24 giờ?
\shortans{2}
\loigiai{
Sự kiện A và C chắc chắn (2 sự kiện).
}
\end{ex}

\begin{ex}
Có bao nhiêu sự kiện 'Không thể' trong 3 sự kiện trên?
\shortans{1}
\loigiai{
Sự kiện B không thể (1 sự kiện).
}
\end{ex}
